\Subhead{5.1.2}{Диаграмма направленности. Коэффициент усиления и коэффициент рассеяния антенны}

Интенсивность излучения $U(\theta,\varphi)=r^2\overline S_r$ измеряется в Вт/ср. Коэффициент направленного действия
\begin{equation}
 D(\theta,\varphi)=\frac{4\pi U}{P_{\rm rad}},\qquad
 G=\eta_{\rm rad}D,\quad \eta_{\rm rad}=\frac{P_{\rm rad}}{P_{\rm accepted}}.
\end{equation}
Усиление учитывает омические и диэлектрические потери, но в этой нормировке не рассогласование фидера. Реализованное усиление $G_{\rm real}=(1-|\Gamma|^2)G$ учитывает рассогласование. При приёме несовпадение поляризаций учитывают отдельным множителем. Для Герца $D_{\max}=3/2$, для идеального тонкого полуволнового вибратора около 1.64. Значения в dBi отсчитываются от изотропного излучателя: $10\log_{10}G$.

\textbf{Коэффициент рассеяния антенны} в антенной радиометрии --- доля мощности вне выбранной области главного лепестка $\Omega_{\rm г}$:
\begin{equation}\label{eq:antennascattering}
 \beta=\frac{\int_{4\pi\setminus\Omega_{\rm г}}U(\Omega)\dd\Omega}
 {\int_{4\pi}U(\Omega)\dd\Omega},\qquad 0\leq\beta\leq1.
\end{equation}
Это безразмерная характеристика бокового/заднего излучения. Для взаимной приёмной антенны та же доля определяет вклад внешнего фона вне главного луча в шум (\Qref{5.1.5}). Определение и связь с температурой: Шарков~\cite{sharkov}, §5.4, с.~184, формулы~(5.38), (5.40).

От неё следует отличать входной коэффициент отражения $\Gamma=(Z_A-Z_0)/(Z_A+Z_0)$ и характеристику переизлучения падающей волны. Последнюю обычно задают эффективной площадью рассеяния
\begin{equation}
 \sigma_{\rm sc}(\widehat{\bm r})=\lim_{r\to\infty}4\pi r^2
 \frac{S_{\rm scattered}(\widehat{\bm r})}{S_{\rm incident}}.
\end{equation}
Она имеет размерность площади и зависит от направлений, поляризаций и нагрузки антенны. Она не тождественна безразмерному усилению или эффективной приёмной площади. Рассеяние содержит структурную часть и вклад тока нагруженного антенного порта; согласование нагрузки не делает всю антенну невидимой.
